\section{Discussion}
\label{sec:discussion}

The preceding results separate three questions that are often conflated in
informal discussions of horizon visibility: whether refraction can produce
occlusion-like ray geometry at all, whether it can reproduce the quantitative
height--distance law of a sphere, and whether the required refractive profile
is physically representative of the atmosphere.  The answer to the first two
questions is affirmative within geometrical optics.  The answer to the third
is negative for ordinary atmospheric stratification.

\subsection{Mathematical possibility versus physical explanation}

The exponential profile provides an explicit constructive example.  In the
paraxial regime it produces rays of constant upward curvature and recovers the
leading spherical horizon law.  Exactly, its optical metric is mapped by a
change of variables to the Euclidean exterior of a circle.  This establishes
that geometric horizon-like occlusion above a flat boundary is not forbidden
by optics as a matter of principle.

That mathematical equivalence is not, by itself, a physical explanation of a
terrestrial observation.  A physical explanation must identify a medium with
the required profile and show that it persists over the relevant horizontal
and vertical scales.  Both the exponential construction and the exact inverse
reconstruction require a refractive index that increases with height.  The
ordinary atmosphere instead has a decreasing refractive index.  The distinction
is therefore one of mechanism, not merely of fitting accuracy.

\subsection{What is reproduced by the construction}

The general inverse formula of Sec.~\ref{sec:inverse-reconstruction} shows that
a prescribed one-sided horizon law $d(h)$ determines a corresponding monotonic
stratified profile.  In particular, the exact spherical law can be reproduced
for all endpoint heights in the idealised model.  Consequently, the following
features can be generated optically above a flat opaque boundary:
\begin{itemize}
    \item a finite limiting distance for objects of fixed height;
    \item disappearance of lower portions before upper portions as distance
          increases;
    \item recovery of previously hidden portions when the observer is raised;
    \item additivity of the observer and target horizon distances.
\end{itemize}
These effects follow from the same grazing-ray structure that determines
geometric tangency on a sphere.

The construction does not automatically reproduce every observable property
of imaging above a sphere.  Equality of the limiting visibility law constrains
the boundary of the visible region, but does not by itself enforce equality of
magnification, angular size, intensity, chromatic dispersion, or image
multiplicity away from that boundary.  The conformal equivalence of the
exponential model is stronger than equality of one horizon law, but even there
the physical height coordinate is related nonlinearly to the radial coordinate
of the circular picture.  Coordinate equivalence must therefore not be confused
with literal identity of the two physical systems.

\subsection{Uniqueness and diagnostic value of the inverse problem}

Within the assumptions of a smooth isotropic index $n(h)$ and a single grazing
branch, the inverse relation
\begin{equation}
    n(h)=n_0\sqrt{1+\frac{1}{[d'(h)]^2}}
\end{equation}
assigns a unique profile once $n_0$ is fixed.  This makes the horizon law a
diagnostic rather than merely a qualitative resemblance.  A proposed optical
mechanism cannot be assessed only by noting that rays bend; its index profile
must generate the measured dependence of limiting distance on height.

The same reasoning also exposes the limits of ad hoc appeals to local
refraction.  A temperature inversion may alter a particular observation, move
an apparent horizon, or create multiple images.  Reproducing one selected ray
or one selected photograph does not establish a universal alternative horizon
geometry.  To reproduce the full family of height-dependent visibility limits,
the refractive structure must satisfy the corresponding inverse profile over
the entire region sampled by those rays.

\subsection{Observational discrimination}

The spherical and refractive models are most difficult to distinguish at small
heights because their leading horizon asymptotics agree.  Discrimination must
therefore rely on information beyond the leading square-root law.  Useful tests
include measurements over a range of observer and target heights, simultaneous
meteorological profiling, repeated observations under changing atmospheric
conditions, and checks for image distortion or multiple ray paths.  A genuine
atmospheric explanation predicts sensitivity to the refractive-index field,
whereas geometric occlusion by a stable surface remains after transient
refractive effects are averaged or independently constrained.

The sign of the required gradient provides an especially direct test.  The
flat-boundary construction that mimics spherical occlusion needs upward bending
rays and hence $dn/dh>0$ in the adopted height coordinate.  Measuring a normal
negative gradient does not merely reduce the strength of the proposed effect;
it implies bending in the opposite direction from that required by the model.

\subsection{Limitations of the analysis}

The analysis is intentionally idealised.  It neglects diffraction, scattering,
absorption, turbulence, wavelength dependence, finite source size, and
three-dimensional horizontal inhomogeneity.  It also assumes a perfectly flat,
opaque lower boundary and a stationary isotropic medium.  These omissions do
not affect the logical distinction between mathematical reconstructibility and
physical plausibility, but they limit direct application to detailed
atmospheric scenes.

Within those limits, the central result is robust: refraction can be engineered
to reproduce horizon-like occlusion above a flat surface, and even to match the
exact spherical height--distance law.  The profile required for that purpose,
however, is a specially constructed increasing-index medium rather than the
ordinary atmosphere.  Refraction is therefore a mathematically valid analogue
of geometric occlusion, but standard atmospheric refraction is not an
interchangeable explanation for it.
